% related.tex
\section{Related Work}

\subsection{Audio Source Separation}

Modern source separation systems decompose polyphonic mixtures into per-instrument stems, enabling downstream analysis of individual instruments.
Spleeter~\cite{hennequin2020spleeter} popularized lightweight U-Net--based separation with pre-trained models for 2-, 4-, and 5-stem splits, while Hybrid Demucs~\cite{defossez2021hybrid} combines waveform and spectrogram branches in a hybrid architecture that achieves state-of-the-art separation quality.
These systems extract isolated audio tracks but do not recover the synthesis parameters that produced them---they answer \emph{what} was played but not \emph{how} the sound was shaped.
Our work takes the output of such separators as input: given an isolated instrument stem, we recover the synthesizer preset that reproduces it.

\subsection{Automatic Music Transcription}

Automatic music transcription (AMT) converts audio into symbolic representations such as MIDI.
MT3~\cite{gardner2022mt3} formulates multi-instrument transcription as a sequence-to-sequence problem using a Transformer encoder--decoder, achieving strong results across diverse timbres.
Basic Pitch provides a lightweight convolutional model for monophonic and polyphonic pitch estimation.
These approaches extract \emph{note-level} information (pitch, onset, duration) but discard the timbral and parametric details of the sound source.
In contrast, our method operates at the \emph{parameter level}: given audio of a single note or phrase, we recover the full set of synthesizer controls---oscillator mix, filter shape, envelope, and effects---that reproduce the timbre.

\subsection{Automatic Synthesizer Programming}

Recovering synthesizer parameters from audio is a long-standing challenge.
InverSynth~\cite{barkan2019inversynth} trains a feed-forward neural network to predict FM synthesis parameters from spectrograms, demonstrating that direct regression can approximate the inverse mapping for a fixed synthesizer.
Sound2Synth~\cite{chen2022sound2synth} extends this to a broader parameter space with a multi-task architecture that jointly predicts discrete and continuous parameters.
DiffMoog~\cite{hayes2024diffmoog} builds a fully differentiable modular synthesizer and optimizes parameters via backpropagation through the signal chain, enabling gradient-based sound matching without a learned model.
SynthRL~\cite{shin2025synthrl} formulates parameter recovery as a reinforcement learning problem, using a policy network to iteratively adjust parameters toward a target timbre.
Yang et al.~\cite{yang2023whitebox} propose a white-box search approach that exploits the known structure of the synthesizer to constrain the search space.

These methods each address different trade-offs between generality, accuracy, and computational cost.
Neural approaches~\cite{barkan2019inversynth,chen2022sound2synth} require large paired datasets for training and are tied to the synthesizer used during data collection.
Differentiable synthesis~\cite{hayes2024diffmoog} avoids dataset dependence but requires the entire signal chain to be differentiable and can become trapped in local minima of the non-convex loss landscape.
Our approach combines the best of both worlds: we use CMA-ES---a derivative-free evolutionary optimizer---for global exploration, followed by gradient-based refinement through our differentiable synthesizer, avoiding both the dataset requirement and the local-minima problem.

\subsection{Differentiable Audio Synthesis}

Differentiable digital signal processing (DDSP)~\cite{engel2020ddsp} introduced the paradigm of embedding classical signal processing components---harmonic oscillators, filtered noise, reverb---as differentiable modules within neural network architectures, enabling end-to-end training with audio reconstruction losses.
FlowSynth~\cite{yang2025flowsynth} extends this idea with normalizing flows to model the distribution of synthesizer parameters, allowing both parameter inference and diverse preset generation.
Our synthesizer draws on the DDSP philosophy: all signal-chain components (oscillators, filters, ADSR envelopes, reverb) are implemented as differentiable PyTorch operations, making the entire rendering pipeline amenable to gradient-based optimization.
However, unlike DDSP-based models that learn an encoder, we optimize parameters directly for each target sound via a hybrid evolutionary--gradient strategy.

\subsection{Evolutionary Optimization for Audio}

The Covariance Matrix Adaptation Evolution Strategy (CMA-ES)~\cite{hansen2016cmaes} is a state-of-the-art derivative-free optimizer for continuous, non-convex problems.
It maintains a multivariate Gaussian search distribution and adapts both its mean and covariance matrix based on the ranking of candidate solutions, providing robust performance on multimodal landscapes without requiring gradient information.
CMA-ES has been applied to sound matching and synthesizer programming tasks where the loss surface is highly non-convex due to the complex, often discontinuous relationship between synthesis parameters and perceptual audio similarity.
Our system uses CMA-ES as the first stage of a two-stage pipeline: CMA-ES locates a promising basin in parameter space, after which Adam gradient descent exploits the differentiability of our synthesizer to refine the solution to high precision.
This hybrid strategy addresses the complementary weaknesses of each approach---CMA-ES provides global exploration but converges slowly near optima, while gradient descent converges quickly but is sensitive to initialization.
